% GENERATED FILE. Edit canonical lessons, then run npm run generate:books.
% canonical-source-hash: fnv1a64:6c86ffc3b6446e94
% canonical-lessons: GE-C358-gesellschaft, GE-C358-demokratie, GE-C358-abstimmen, GE-C358-abstimmung, GE-C358-beitragen

\chapter{Die Gesellschaft}
\label{ch:ge-die-gesellschaft}
\hlchaptermodality{358}

\begin{chapteropening}
\textbf{By the end of this chapter:} \emph{I can talk about society and democracy, voting, and what people contribute.}

Complete the last lesson of chapter 358: I can talk about society and democracy, voting, and what people contribute.
\end{chapteropening}

\section[die Gesellschaft]{die Gesellschaft — society; company; the firm}

\label{lesson:GE-C358-gesellschaft}



\begin{quote}
Before the new one: say the German for to support, then the German for the support.
\end{quote}

\subsection*{You'll want to know: die Gesellschaft}
\textbf{die Gesellschaft} — \textquotedblleft{}society; company; the firm\textquotedblright{}.

\begin{quote}
In unserer Gesellschaft leben immer mehr alte Menschen.
\end{quote}

\begin{quote}
\emph{More and more old people live in our society.}
\end{quote}

Plural: \emph{die Gesellschaften}. Also company, being with others: \emph{in guter Gesellschaft}; and a firm, as in \emph{GmbH}.

\begin{cousinweb}
\emph{Gesellschaft} is inherited from Middle High German \emph{gesellescaft} and Old High German \emph{giselliscaft} (EtymDB), built on \emph{Geselle} \textquotedblleft{}companion\textquotedblright{}, usually explained as first someone who shared your hall, with \emph{-schaft}.
\end{cousinweb}

\begin{grammarlens}[title={What you've built}]
One more word for telling, explaining and describing at B1.
\end{grammarlens}

\subsection*{Your turn}
\begin{itemize}
\raggedright
  \item \emph{Say it:} \emph{die Gesellschaft}
  \item \emph{Say it:} \emph{die Gesellschaft}, once more
  \item \emph{Say it:} a sentence of your own with \emph{die Gesellschaft}
  \item \emph{Recall:} say the German for to support, then the German for the support, then say \emph{die Gesellschaft} again
\end{itemize}

\begin{tcolorbox}[breakable,colback=teal!4,colframe=teal!35!black,title={Before you move on}]
What does die Gesellschaft mean? (\textquotedblleft{}Society; company; the firm\textquotedblright{}.) Say it once more.
\end{tcolorbox}
\section[die Demokratie]{die Demokratie — democracy}

\label{lesson:GE-C358-demokratie}



\begin{quote}
Before the new one: say the German for the support, then the German for society.
\end{quote}

\subsection*{You'll want to know: die Demokratie}
\textbf{die Demokratie} — \textquotedblleft{}democracy\textquotedblright{}.

\begin{quote}
In einer Demokratie wählen die Bürger das Parlament.
\end{quote}

\begin{quote}
\emph{In a democracy the citizens elect parliament.}
\end{quote}

Plural: \emph{die Demokratien}. The stress is on the last syllable: \emph{de-mo-kra-TIE}.

\begin{cousinweb}
\emph{Demokratie} was borrowed from French \emph{démocratie} (EtymDB), from Latin \emph{dēmocratia} and Greek \emph{dēmokratía}, built from \emph{dêmos} \textquotedblleft{}the common people\textquotedblright{} and \emph{krátos} \textquotedblleft{}strength, power\textquotedblright{}.
\end{cousinweb}

\begin{grammarlens}[title={What you've built}]
One more word for telling, explaining and describing at B1.
\end{grammarlens}

\subsection*{Your turn}
\begin{itemize}
\raggedright
  \item \emph{Say it:} \emph{die Demokratie}
  \item \emph{Say it:} \emph{die Demokratie}, once more
  \item \emph{Say it:} a sentence of your own with \emph{die Demokratie}
  \item \emph{Recall:} say the German for the support, then the German for society, then say \emph{die Demokratie} again
\end{itemize}

\begin{tcolorbox}[breakable,colback=teal!4,colframe=teal!35!black,title={Before you move on}]
What does die Demokratie mean? (\textquotedblleft{}Democracy\textquotedblright{}.) Say it once more.
\end{tcolorbox}
\section[abstimmen]{abstimmen — to vote (on); to coordinate}

\label{lesson:GE-C358-abstimmen}



\begin{quote}
Before the new one: say the German for society, then the German for democracy.
\end{quote}

\subsection*{You'll want to know: abstimmen}
\textbf{abstimmen} — \textquotedblleft{}to vote (on); to coordinate\textquotedblright{}.

\begin{quote}
Am Sonntag stimmen die Bürger über den neuen Bahnhof ab.
\end{quote}

\begin{quote}
\emph{On Sunday the citizens vote on the new station.}
\end{quote}

Separable: \emph{sie stimmen ab}, \emph{sie haben abgestimmt}. \emph{Abstimmen über} + accusative. \emph{Termine abstimmen} — to coordinate dates.

\begin{cousinweb}
\emph{abstimmen} is built from \emph{ab} and \emph{stimmen}, a verb made from \emph{Stimme} \textquotedblleft{}voice, vote\textquotedblright{} (EtymDB): you give your vote.
\end{cousinweb}

\begin{grammarlens}[title={What you've built}]
One more word for telling, explaining and describing at B1.
\end{grammarlens}

\subsection*{Your turn}
\begin{itemize}
\raggedright
  \item \emph{Say it:} \emph{abstimmen}
  \item \emph{Say it:} \emph{abstimmen}, once more
  \item \emph{Say it:} a sentence of your own with \emph{abstimmen}
  \item \emph{Recall:} say the German for society, then the German for democracy, then say \emph{abstimmen} again
\end{itemize}

\begin{tcolorbox}[breakable,colback=teal!4,colframe=teal!35!black,title={Before you move on}]
What does abstimmen mean? (\textquotedblleft{}To vote (on); to coordinate\textquotedblright{}.) Say it once more.
\end{tcolorbox}
\section[die Abstimmung]{die Abstimmung — the vote, the ballot}

\label{lesson:GE-C358-abstimmung}



\begin{quote}
Before the new one: say the German for democracy, then the German for to vote.
\end{quote}

\subsection*{You'll want to know: die Abstimmung}
\textbf{die Abstimmung} — \textquotedblleft{}the vote, the ballot\textquotedblright{}.

\begin{quote}
Das Ergebnis der Abstimmung war sehr knapp.
\end{quote}

\begin{quote}
\emph{The result of the vote was very close.}
\end{quote}

Plural: \emph{die Abstimmungen}. In Switzerland \emph{die Volksabstimmung}, a vote of the whole people, takes place several times a year.

\begin{cousinweb}
\emph{Abstimmung} is the noun of \emph{abstimmen}, formed with \emph{-ung}; the verb is made from \emph{Stimme} \textquotedblleft{}voice, vote\textquotedblright{} (EtymDB).
\end{cousinweb}

\begin{grammarlens}[title={What you've built}]
One more word for telling, explaining and describing at B1.
\end{grammarlens}

\subsection*{Your turn}
\begin{itemize}
\raggedright
  \item \emph{Say it:} \emph{die Abstimmung}
  \item \emph{Say it:} \emph{die Abstimmung}, once more
  \item \emph{Say it:} a sentence of your own with \emph{die Abstimmung}
  \item \emph{Recall:} say the German for democracy, then the German for to vote, then say \emph{die Abstimmung} again
\end{itemize}

\begin{tcolorbox}[breakable,colback=teal!4,colframe=teal!35!black,title={Before you move on}]
What does die Abstimmung mean? (\textquotedblleft{}The vote, the ballot\textquotedblright{}.) Say it once more.
\end{tcolorbox}
\section[beitragen]{beitragen — to contribute}

\label{lesson:GE-C358-beitragen}



\begin{quote}
Before the new one: say the German for to vote, then the German for the vote.
\end{quote}

\subsection*{You'll want to know: beitragen}
\textbf{beitragen} — \textquotedblleft{}to contribute\textquotedblright{}.

\begin{quote}
Jeder kann etwas zu einem schönen Fest beitragen.
\end{quote}

\begin{quote}
\emph{Everyone can contribute something to a nice party.}
\end{quote}

Separable and strong: \emph{er trägt bei}, \emph{er hat beigetragen}. \emph{Der Beitrag} — the contribution; also a fee or an article.

\begin{cousinweb}
\emph{beitragen} is built from \emph{bei} \textquotedblleft{}by, to\textquotedblright{} and \emph{tragen} \textquotedblleft{}to carry\textquotedblright{}, from Proto-Germanic \emph{draganą} (EtymDB): you carry your share to the pile.
\end{cousinweb}

\begin{grammarlens}[title={What you've built}]
One more word for telling, explaining and describing at B1.
\end{grammarlens}

\subsection*{Your turn}
\begin{itemize}
\raggedright
  \item \emph{Say it:} \emph{beitragen}
  \item \emph{Say it:} \emph{beitragen}, once more
  \item \emph{Say it:} a sentence of your own with \emph{beitragen}
  \item \emph{Recall:} say the German for to vote, then the German for the vote, then say \emph{beitragen} again
\end{itemize}

\begin{tcolorbox}[breakable,colback=teal!4,colframe=teal!35!black,title={Before you move on}]
What does beitragen mean? (\textquotedblleft{}To contribute\textquotedblright{}.) Say it once more.
\end{tcolorbox}
